\documentclass[12pt,a4paper]{ctexart}
\usepackage[left=2.5cm,right=2.5cm,top=2.3cm,bottom=2.3cm]{geometry}
\usepackage{xcolor}
\usepackage{tcolorbox}
\usepackage{titlesec}
\usepackage{fancyhdr}
\usepackage{enumitem}
\usepackage{tabularx}
\usepackage{amsmath,amssymb}
\usepackage{hyperref}
\tcbuselibrary{skins,breakable}

\definecolor{MainColor}{HTML}{2A6F73}
\definecolor{SecondColor}{HTML}{3CAEA3}
\definecolor{LightMain}{HTML}{EAF7F6}
\definecolor{TitleDark}{HTML}{1F3A40}
\definecolor{TextGray}{HTML}{555555}
\definecolor{BorderGray}{HTML}{CBD5D6}
\definecolor{WarnColor}{HTML}{F2B35D}
\definecolor{ErrorColor}{HTML}{C94C4C}
\definecolor{DefineColor}{HTML}{6C63B7}

\linespread{1.25}
\setlength{\parindent}{2em}
\setlength{\parskip}{0.25em}
\hypersetup{colorlinks=true,linkcolor=MainColor,urlcolor=MainColor}
\pagestyle{fancy}
\fancyhf{}
\fancyfoot[L]{\textcolor{MainColor}{机器学习笔记}}
\fancyfoot[R]{\textcolor{MainColor}{\thepage}}
\renewcommand{\headrulewidth}{0pt}
\renewcommand{\footrulewidth}{0pt}

\newcommand{\topbar}[2]{\noindent{\small\textcolor{MainColor}{#1}}\hfill{\small\bfseries\textcolor{TitleDark}{#2}}\par\vspace{0.35em}{\color{BorderGray}\hrule}\vspace{1.2em}}
\newcommand{\notetitlecard}[6]{
\begin{tcolorbox}[enhanced,colback=white,colframe=BorderGray,boxrule=0.6pt,arc=2mm,left=12pt,right=12pt,top=10pt,bottom=10pt,borderline west={4pt}{0pt}{SecondColor}]
\begin{tabularx}{\linewidth}{X r}
{\small\textcolor{TextGray}{#1}} & {\small\textcolor{TextGray}{#2}} \\[0.9em]
\multicolumn{2}{l}{{\Huge\bfseries\textcolor{MainColor}{#3}}} \\[0.45em]
\multicolumn{2}{l}{{\large\textcolor{SecondColor!90!black}{#4}}} \\[0.45em]
\multicolumn{2}{l}{{\small\textcolor{TextGray}{课程：#5 \qquad 作者：#6}}}
\end{tabularx}\end{tcolorbox}\vspace{1em}}
\newtcolorbox{summarybox}[1][本节摘要]{enhanced,breakable,colback=LightMain,colframe=SecondColor,colbacktitle=SecondColor,coltitle=white,title={#1},fonttitle=\bfseries,boxrule=0.6pt,arc=2mm,left=12pt,right=12pt,top=10pt,bottom=10pt}
\newtcolorbox{definitionbox}[1][定义]{enhanced,breakable,colback=DefineColor!6,colframe=DefineColor,colbacktitle=DefineColor,coltitle=white,title={#1},fonttitle=\bfseries,boxrule=0.6pt,arc=2mm,left=12pt,right=12pt,top=10pt,bottom=10pt}
\newtcolorbox{keybox}[1][重点]{enhanced,breakable,colback=MainColor!5,colframe=MainColor,colbacktitle=MainColor,coltitle=white,title={#1},fonttitle=\bfseries,boxrule=0.6pt,arc=2mm,left=12pt,right=12pt,top=10pt,bottom=10pt}
\newtcolorbox{examplebox}[1][例题]{enhanced,breakable,colback=white,colframe=SecondColor,colbacktitle=SecondColor!85!black,coltitle=white,title={#1},fonttitle=\bfseries,boxrule=0.6pt,arc=2mm,left=12pt,right=12pt,top=10pt,bottom=10pt}
\newtcolorbox{mistakebox}[1][易错点]{enhanced,breakable,colback=ErrorColor!5,colframe=ErrorColor,colbacktitle=ErrorColor,coltitle=white,title={#1},fonttitle=\bfseries,boxrule=0.6pt,arc=2mm,left=12pt,right=12pt,top=10pt,bottom=10pt}
\newtcolorbox{tipbox}[1][提示]{enhanced,breakable,colback=WarnColor!10,colframe=WarnColor,colbacktitle=WarnColor,coltitle=white,title={#1},fonttitle=\bfseries,boxrule=0.6pt,arc=2mm,left=12pt,right=12pt,top=10pt,bottom=10pt}
\titleformat{\section}{\Large\bfseries\color{MainColor}}{\thesection}{0.8em}{}
\titleformat{\subsection}{\large\bfseries\color{TitleDark}}{\thesubsection}{0.8em}{}
\titleformat{\subsubsection}{\normalsize\bfseries\color{SecondColor!80!black}}{\thesubsubsection}{0.8em}{}
\setlist[itemize]{leftmargin=2.2em,itemsep=0.3em,topsep=0.4em}
\setlist[enumerate]{leftmargin=2.2em,itemsep=0.3em,topsep=0.4em}
\newcommand{\important}[1]{\textbf{\textcolor{MainColor}{#1}}}
\newcommand{\warning}[1]{\textbf{\textcolor{ErrorColor}{#1}}}
\newcommand{\concept}[1]{\textbf{\textcolor{DefineColor}{#1}}}

\begin{document}
\topbar{吴恩达机器学习}{学习笔记}
\notetitlecard{Study Notes Series}{2026 Spring}{吴恩达机器学习笔记(7)}{}{机器学习}{C++与草稿纸}

\begin{summarybox}
本周介绍支持向量机（SVM）：用带间隔的优化目标寻找稳健的分类边界，理解大间隔分类的几何直觉，并使用核函数构造非线性分类器。重点是目标函数、参数 $C$、核参数和实际选型。
\end{summarybox}

\section{支持向量机的优化目标}
\begin{definitionbox}[从逻辑回归到 SVM]
逻辑回归的单样本损失可以看作两段惩罚函数。SVM 用近似折线损失替代对数损失：正类希望 $\theta^Tx\geq1$，负类希望 $\theta^Tx\leq-1$。设 $y\in\{-1,1\}$，则约束可写作
\[
y\,\theta^Tx\geq1.
\]
\end{definitionbox}
\begin{keybox}[硬间隔与软间隔目标]
线性 SVM 的简化目标是
\[
\underset{\theta}{\operatorname{minimize}}
\;\frac12\sum_{j=1}^{n}\theta_j^2+C\sum_{i=1}^{m}\mathrm{cost}_1(\theta^Tx^{(i)},y^{(i)}),
\]
其中 $\mathrm{cost}_1$ 和 $\mathrm{cost}_0$ 分别对应正类和负类的折线损失。第一项倾向于让参数范数小，第二项惩罚分类错误或落入间隔内的样本。实际实现通常使用凸二次规划求解，而不是手写梯度下降。
\end{keybox}

\section{大间隔分类的直观理解}
\begin{definitionbox}[几何间隔]
分类边界由 $\theta^Tx=0$ 给出。参数向量 $\theta$ 垂直于边界，点到边界的距离与 $|\theta^Tx|/\|\theta\|$ 成正比。SVM 不仅要求分类正确，还希望正负样本离边界有尽可能大的安全间隔。
\end{definitionbox}
\begin{examplebox}[为什么间隔更稳健]
若存在多个能把训练样本分开的直线，靠近样本的边界对噪声和新样本更敏感；大间隔边界尽量远离两侧样本，通常具有更好的泛化能力。落在间隔边界附近、决定最终分界的样本称为支持向量。
\end{examplebox}
\begin{tipbox}[$C$ 的作用]
$C$ 是误分类惩罚的权重：$C$ 很大时更强调训练集拟合，间隔可能变窄并提高方差；$C$ 较小时允许少量训练错误，间隔更宽，通常更抗噪声。实际应使用交叉验证选择 $C$。
\end{tipbox}

\section{核函数与非线性边界}
\begin{definitionbox}[特征映射]
对输入 $x$ 构造一组相似度特征 $f_1(x),\ldots,f_m(x)$，再在这些新特征上训练线性 SVM，就可以得到原始空间中的非线性边界。核函数直接计算样本相似度，避免显式构造高维特征。
\end{definitionbox}
\begin{examplebox}[高斯核]
高斯核（RBF kernel）常写为
\[
K(x,l)=\exp\!\left(-\frac{\|x-l\|^2}{2\sigma^2}\right).
\]
当 $x$ 接近地标 $l$ 时，核值接近 $1$；距离较远时接近 $0$。$\sigma$ 控制相似度衰减速度：$\sigma$ 小会形成局部、复杂的边界，$\sigma$ 大则边界更平滑。
\end{examplebox}
\begin{mistakebox}[核函数使用条件]
\begin{itemize}
  \item 特征尺度差异很大时，距离会被大尺度特征主导，应先进行特征缩放；
  \item 不是所有相似度都能作为合法核函数，常见实现应使用经过验证的核；
  \item $C$ 和 $\sigma$ 都是超参数，应在交叉验证集上联合选择；
  \item 线性核在特征已经足够多或样本极大时可能更快，不必默认使用高斯核。
\end{itemize}
\end{mistakebox}

\section{如何使用 SVM}
\begin{keybox}[模型选择建议]
样本数 $m$ 很大而特征数 $n$ 较小，可考虑线性 SVM；$m$ 中等时可使用高斯核；$n$ 和 $m$ 都很大时，线性模型或近似核方法通常更实际。神经网络也能得到类似的非线性分类能力，但训练和部署成本不同。
\end{keybox}
\begin{tipbox}[多类别 SVM]
多类别任务可以训练多个一对多分类器，选择决策函数值最大的类别；也可以使用库提供的多类策略。预测分数不一定是校准概率，需要概率解释时应额外进行校准。
\end{tipbox}
\section{本周知识脉络}
\begin{keybox}[大间隔与核技巧]
SVM 以参数范数和分类损失共同定义优化目标；几何上寻找离样本更远的边界；核函数把相似度变成可学习的高维表示；$C$ 控制误差惩罚，$\sigma$ 控制高斯核的局部程度。数据缩放、交叉验证和计算规模决定最终选型。
\end{keybox}
\end{document}
